%=====================================================================
\chapter{Describing Data: Centre, Spread and Shape}\label{ch:descriptive}
%=====================================================================
\locator{1}{Part I --- Foundations of Data}

\begin{outcomes}
By the end of this chapter you will be able to:
\begin{tight}
  \item \textbf{Compute} and \textbf{choose between} the mean, median and mode for a given variable.
  \item \textbf{Compute} range, variance, standard deviation and the interquartile range, and say what each is robust to.
  \item \textbf{Read} a boxplot and identify outliers by the $1.5 \times \text{IQR}$ rule.
  \item \textbf{Recognise} skew from a histogram and predict how it displaces the mean relative to the median.
  \item \textbf{Identify} which common distribution a variable plausibly follows and why that matters.
\end{tight}
\end{outcomes}

\begin{prereq}
\chref{data} (measurement scales --- they decide which summary is legal) and
\chref{math} (summation notation).
\end{prereq}

\begin{keyterms}
mean $\bullet$ median $\bullet$ mode $\bullet$ range $\bullet$ variance $\bullet$
standard deviation $\bullet$ quartile $\bullet$ interquartile range $\bullet$
outlier $\bullet$ skewness $\bullet$ normal, uniform, binomial, Poisson and
long-tailed distributions $\bullet$ robustness
\end{keyterms}

\section{Why Summarise at All}

You cannot look at 400{,}000 rows. A summary replaces a dataset with a handful of
numbers that preserve what matters and discard what does not --- and the entire
skill lies in knowing which is which. Every summary is a deliberate loss of
information.

\section{Measures of Centre}

\begin{definitionbox}[Mean, Median, Mode]
The \term{mean} $\bar{x} = \frac{1}{n}\sum_{i=1}^{n} x_i$ is the arithmetic
average.\\[2pt]
The \term{median} is the middle value when the data is sorted (the average of the
two middle values if $n$ is even).\\[2pt]
The \term{mode} is the most frequently occurring value --- the only centre that
is meaningful for nominal data.
\end{definitionbox}

\begin{worked}[When the Mean Lies]
Response times in milliseconds for seven API calls:
$$12,\; 14,\; 15,\; 15,\; 17,\; 19,\; 8400$$

\textbf{Mean:} $(12+14+15+15+17+19+8400)/7 = 8492/7 \approx 1213$~ms.

\textbf{Median:} sorted, the 4th of 7 values is $15$~ms.

\textbf{Mode:} $15$~ms.

\smallskip
\textbf{Interpretation.} The mean of 1213~ms describes not a single one of the
seven calls. Six were under 20~ms; one timed out. Reporting ``average response
time is 1.2 seconds'' would be technically correct and completely misleading.

\textbf{The rule this illustrates:} the mean is \term{sensitive} to extreme
values, the median is \term{robust} to them. Report the median for anything with a
long tail --- response times, salaries, session durations, file sizes --- which in
practice is most operational data you will ever meet.
\end{worked}

\rowcolors{2}{white}{lightbg}
\begin{longtable}{@{}p{2.0cm}p{3.0cm}p{4.0cm}p{5.2cm}@{}}
\toprule
\rowcolor{primary!15}
\textbf{Measure} & \textbf{Valid for} & \textbf{Robust to outliers?} & \textbf{Use when} \\
\midrule
\endhead
Mean & Interval, ratio & No --- one value can move it arbitrarily & Data is roughly symmetric and you need to do further arithmetic \\
Median & Ordinal, interval, ratio & Yes --- unaffected by how extreme the extremes are & Data is skewed, or outliers are present and genuine \\
Mode & Any, including nominal & Yes & The variable is categorical, or you want the most typical case \\
\bottomrule
\end{longtable}

\section{Measures of Spread}

Two datasets can share a mean and be nothing alike. Spread is what distinguishes
``every server responds in about 40~ms'' from ``servers respond in 5~ms or 200~ms
depending on the day''.

\begin{definitionbox}[Variance and Standard Deviation]
The \term{sample variance} is
\[
s^2 \;=\; \frac{1}{n-1}\sum_{i=1}^{n}\left(x_i - \bar{x}\right)^2
\]
and the \term{standard deviation} is $s = \sqrt{s^2}$. The standard deviation is
the more useful of the two because it is expressed in the \emph{same units} as the
data.
\end{definitionbox}

\begin{conceptbox}[Why Divide by $n-1$?]
Because $\bar{x}$ was itself estimated from the same data, the deviations are very
slightly too small on average. Dividing by $n-1$ rather than $n$ corrects the
bias. For large $n$ the difference is negligible; for the small samples you will
meet in project management data it is not.
\end{conceptbox}

\begin{definitionbox}[Quartiles and the Interquartile Range]
Sort the data. \term{Q1} is the value below which 25\% falls, \term{Q2} is the
median, \term{Q3} is the value below which 75\% falls. The \term{interquartile
range} is $\text{IQR} = Q3 - Q1$: the width of the middle half of the data. Like
the median, it is robust.
\end{definitionbox}

\dsfig{%
\begin{tikzpicture}[font=\scriptsize]
% ---- the box ----
\def\qone{2.4} \def\med{4.0} \def\qthree{6.2}
\def\wlo{0.8}  \def\whi{8.5}
\draw[fill=secondary!12, draw=secondary, line width=1.1pt] (\qone,-0.55) rectangle (\qthree,0.55);
\draw[primary, line width=1.8pt] (\med,-0.55) -- (\med,0.55);
\draw[secondary, line width=1.1pt] (\qone,0) -- (\wlo,0);
\draw[secondary, line width=1.1pt] (\qthree,0) -- (\whi,0);
\draw[secondary, line width=1.1pt] (\wlo,-0.3) -- (\wlo,0.3);
\draw[secondary, line width=1.1pt] (\whi,-0.3) -- (\whi,0.3);
% outliers
\foreach \x in {10.1, 10.9}
  \fill[accent] (\x,0) circle (2.2pt);

% ---- labels ----
\node[text=primary, font=\scriptsize\bfseries] at (\med,0.95) {median};
\node[text=secondary!60!black] at (\qone,-0.95) {Q1};
\node[text=secondary!60!black] at (\qthree,-0.95) {Q3};
\node[text=secondary!60!black, align=center] at (\wlo,-1.05) {lower\\whisker};
\node[text=secondary!60!black, align=center] at (\whi,-1.05) {upper\\whisker};
\node[text=accent, align=center, font=\scriptsize\bfseries] at (10.5,0.95) {outliers};

% IQR brace
\draw[decorate, decoration={brace, amplitude=4pt}, gray!65]
     (\qone,0.68) -- (\qthree,0.68)
     node[midway, above=4pt, text=gray!55!black] {IQR $=Q3-Q1$: the middle 50\%};

% whisker rule
\draw[decorate, decoration={brace, amplitude=4pt, mirror}, gray!65]
     (\qthree,-1.55) -- (\whi,-1.55)
     node[midway, below=4pt, text=gray!55!black] {up to $1.5\times$IQR};

\draw[-{Stealth[length=2mm]}, gray!55] (-0.2,-2.15) -- (11.4,-2.15);
\node[text=gray!55!black] at (11.4,-2.5) {value};
\end{tikzpicture}}%
{Anatomy of a boxplot. The box holds the middle 50\% of the data, the whiskers reach the
furthest point within $1.5\times$IQR of the box, and anything beyond is drawn individually
as a candidate outlier. A boxplot shows centre, spread, skew and outliers in one
object.}{boxplot-anatomy}

\begin{definitionbox}[The $1.5 \times$ IQR Rule]
A value is flagged as a candidate \term{outlier} if it lies below
$Q1 - 1.5\,\text{IQR}$ or above $Q3 + 1.5\,\text{IQR}$. ``Candidate'' is the
operative word: the rule identifies unusual values, not wrong ones.
\end{definitionbox}

\begin{worked}[Applying the Outlier Rule]
Login durations (minutes) for one account over ten sessions:
$$4,\;5,\;6,\;6,\;7,\;8,\;9,\;11,\;14,\;96$$

\textbf{Q1} (25th percentile) $= 6$. \quad \textbf{Q3} (75th percentile) $= 11$.

\textbf{IQR} $= 11 - 6 = 5$.

\textbf{Fences:} lower $= 6 - 1.5(5) = -1.5$; upper $= 11 + 1.5(5) = 18.5$.

\textbf{Flagged:} 96 only.

\smallskip
\textbf{And now the part that matters.} What you do next depends entirely on
domain, not statistics:
\begin{tight}
  \item \emph{Learning analytics:} probably a student who left a tab open. Cap or
        remove.
  \item \emph{Cybersecurity:} a 96-minute session on an account that normally runs
        6 minutes is precisely what you are paid to notice. \textbf{Do not
        remove it.} In \chref{unsupervised} the outlier becomes the target.
\end{tight}
Never delete outliers because a rule flagged them. Investigate, decide, and record
the decision.
\end{worked}

\section{Shape: Skew and the Distributions You Will Meet}

\dsfig{%
\begin{tikzpicture}
\begin{groupplot}[
  group style={group size=3 by 1, horizontal sep=0.85cm},
  width=5.5cm, height=4.0cm,
  axis lines=left, ymin=0, ytick=\empty,
  tick label style={font=\tiny}, title style={font=\scriptsize\bfseries},
  xlabel style={font=\tiny}, enlarge x limits=0.02,
]
% --- right skew
\nextgroupplot[title={\color{accent}Right-skewed (positive)}, xmin=0, xmax=8, xlabel={value}]
\addplot[draw=accent!70, fill=accent!14, line width=0.9pt, domain=0.05:8, samples=140]
  {2.6*exp(-(ln(x)-0.55)^2/(2*0.55^2))/(x*0.55*2.5066)};
\draw[greenhl, line width=1pt, dashed] (axis cs:1.73,0) -- (axis cs:1.73,1.05);
\draw[primary, line width=1pt] (axis cs:2.02,0) -- (axis cs:2.02,1.05);
\node[font=\tiny, greenhl!60!black, anchor=east] at (axis cs:1.65,1.18) {median};
\node[font=\tiny, primary, anchor=west] at (axis cs:2.1,1.18) {mean};
\node[font=\tiny, text=gray!55!black, align=center] at (axis cs:5.6,0.72)
     {tail pulls\\the mean right};

% --- symmetric
\nextgroupplot[title={\color{greenhl}Symmetric}, xmin=-4, xmax=4, xlabel={value}]
\addplot[draw=greenhl!70, fill=greenhl!14, line width=0.9pt, domain=-4:4, samples=140]
  {exp(-x^2/2)/2.5066};
\draw[primary, line width=1.2pt] (axis cs:0,0) -- (axis cs:0,0.42);
\node[font=\tiny, primary, align=center, anchor=south] at (axis cs:0,0.43)
     {mean = median\\= mode};

% --- left skew
\nextgroupplot[title={\color{secondary}Left-skewed (negative)}, xmin=0, xmax=8, xlabel={value}]
\addplot[draw=secondary!70, fill=secondary!14, line width=0.9pt, domain=0.05:7.95, samples=140]
  {2.6*exp(-(ln(8-x)-0.55)^2/(2*0.55^2))/((8-x)*0.55*2.5066)};
\draw[greenhl, line width=1pt, dashed] (axis cs:6.27,0) -- (axis cs:6.27,1.05);
\draw[primary, line width=1pt] (axis cs:5.98,0) -- (axis cs:5.98,1.05);
\node[font=\tiny, primary, anchor=east] at (axis cs:5.9,1.18) {mean};
\node[font=\tiny, greenhl!60!black, anchor=west] at (axis cs:6.35,1.18) {median};
\end{groupplot}
\end{tikzpicture}}%
{Skew and what it does to the centre. The mean is dragged towards the long tail; the
median stays put. If mean and median differ substantially, the distribution is skewed and
you should report the median.}{skewness}

\begin{conceptbox}[A Diagnostic You Can Do in Your Head]
Compare the mean and the median.
If $\bar{x} \gg \text{median}$, the data is right-skewed --- there is a long tail
of large values. If $\bar{x} \ll \text{median}$, it is left-skewed. If they are
close, it is roughly symmetric. This one comparison tells you more about a variable
than any single number.
\end{conceptbox}

\subsection{Four Distributions Worth Recognising}

\dsfig{%
\begin{tikzpicture}
\begin{groupplot}[
  group style={group size=4 by 1, horizontal sep=0.65cm},
  width=4.4cm, height=3.5cm,
  axis lines=left, ymin=0, ytick=\empty,
  tick label style={font=\tiny}, title style={font=\tiny\bfseries},
]
\nextgroupplot[title={Normal}, xmin=-4, xmax=4]
\addplot[draw=primary, fill=primary!15, line width=0.9pt, domain=-4:4, samples=100]
  {exp(-x^2/2)/2.5066};

\nextgroupplot[title={Uniform}, xmin=-0.5, xmax=4.5, ymax=0.35]
\addplot[draw=secondary, fill=secondary!15, line width=0.9pt] coordinates
  {(0,0) (0,0.25) (4,0.25) (4,0)};

\nextgroupplot[title={Binomial ($n{=}10$, $p{=}0.3$)}, xmin=-0.6, xmax=9, ybar,
               bar width=4pt, ymax=0.30]
\addplot[draw=greenhl, fill=greenhl!30] coordinates
  {(0,0.028) (1,0.121) (2,0.233) (3,0.267) (4,0.200) (5,0.103) (6,0.037) (7,0.009) (8,0.001)};

\nextgroupplot[title={Long-tailed}, xmin=0, xmax=8]
\addplot[draw=accent, fill=accent!15, line width=0.9pt, domain=0.05:8, samples=140]
  {2.6*exp(-(ln(x)-0.35)^2/(2*0.75^2))/(x*0.75*2.5066)};
\end{groupplot}
\end{tikzpicture}}%
{Four shapes you will meet constantly. Recognising which one a variable follows tells you
which summary to report, which model assumptions are safe, and whether ``three standard
deviations'' means anything.}{four-distributions}

\rowcolors{2}{white}{defbg}
\begin{longtable}{@{}p{2.4cm}p{5.0cm}p{7.0cm}@{}}
\toprule
\rowcolor{primary!15}
\textbf{Distribution} & \textbf{Arises when} & \textbf{You will meet it as} \\
\midrule
\endhead
Normal & Many small independent effects add up & Exam marks, measurement error, sensor noise around a set point \\
Uniform & Every value equally likely & Random IDs, hash outputs, a well-shuffled sample \\
Binomial & $n$ independent yes/no trials with fixed probability & Number of quizzes passed out of 10; number of packets dropped out of 1000 \\
Poisson & Counts of rare events in a fixed window & Failed logins per minute; defects per release; sensor faults per day \\
Long-tailed & Multiplicative effects, ``rich get richer'' & Response times, file sizes, session lengths, incident cost \\
\bottomrule
\end{longtable}

\begin{alertbox}[Not Everything Is Normal]
The normal distribution is the default assumption of a great deal of classical
statistics, and a great deal of operational data violates it badly. Response times
cannot be negative and have a long right tail; counts of rare security events are
Poisson, not normal. Applying a rule like ``flag anything beyond three standard
deviations'' to a long-tailed variable will flag a large fraction of perfectly
normal traffic. \chref{inference} returns to when normality can be assumed
anyway.
\end{alertbox}

%---------------------------------------------------------------------
\begin{fourlenses}{Summarising a Variable}
\lensLA{Time spent on an assignment is strongly right-skewed: most students take
2--4 hours, a few take thirty. Report the \emph{median} hours, not the mean, or
every published figure will overstate typical effort and demoralise students who
compare themselves to it. The IQR is the honest measure of ``how much do students
vary''.}

\lensPM{Task completion times are right-skewed for the same reason: a task can
overrun without limit but cannot finish in less than zero. This is why
median-based forecasting beats mean-based estimation, and why experienced
estimators quote ranges (Q1--Q3) rather than points. A sprint planned on mean
velocity will miss more often than it hits.}

\lensIOT{Sensor readings around a controlled set point are approximately normal,
so the standard deviation is meaningful and a $3\sigma$ rule is defensible.
Battery life and time-between-failures are long-tailed, so it is not. The same
fleet therefore needs different summaries for different channels --- do not apply
one rule across the telemetry.}

\lensSEC{Failed logins per minute per account is a Poisson count, usually with a
mode of zero. The mean is a poor description and the standard deviation is nearly
meaningless. What matters is the \emph{tail}: the 99.9th percentile, and how far
above it a given account sits. Note also that the outlier you would delete in the
learning-analytics lens is the alert you would escalate here.}
\end{fourlenses}

%---------------------------------------------------------------------
\begin{lab}[Centre, Spread and Shape in Eight Lines]
\begin{lstlisting}[language=Python]
import numpy as np, pandas as pd

rt = pd.Series([12, 14, 15, 15, 17, 19, 8400])   # response times (ms)

print("mean   :", rt.mean())        # 1213.1  <- describes nothing
print("median :", rt.median())      # 15.0    <- describes six of seven
print("std    :", rt.std())         # huge, and equally misleading
print("IQR    :", rt.quantile(.75) - rt.quantile(.25))

q1, q3 = rt.quantile(.25), rt.quantile(.75)
iqr = q3 - q1
outliers = rt[(rt < q1 - 1.5*iqr) | (rt > q3 + 1.5*iqr)]
print("flagged:", list(outliers))

# The skew diagnostic you can do in your head:
print("mean >> median ?", rt.mean() > 2 * rt.median())   # True => right-skewed
\end{lstlisting}

\textbf{R equivalent} (the repo's existing Quarto labs use this style):
\begin{lstlisting}[language=R]
rt <- c(12, 14, 15, 15, 17, 19, 8400)
summary(rt)          # min, Q1, median, mean, Q3, max in one call
IQR(rt)
boxplot(rt, horizontal = TRUE)
\end{lstlisting}
\end{lab}

%---------------------------------------------------------------------
\begin{checkpoint}
\begin{tightnum}
  \item A variable has mean 45 and median 12. What shape is it, and which figure
        would you publish?
  \item For the data $2, 4, 4, 5, 9$: compute the mean, the median and the sample
        standard deviation.
  \item Why is the mode the only measure of centre available for the variable
        \texttt{protocol} $\in$ \{TCP, UDP, ICMP\}?
\end{tightnum}
\end{checkpoint}

%---------------------------------------------------------------------
\begin{chaptersummary}
\begin{tight}
  \item The mean is sensitive to extremes; the median and IQR are robust. Skewed
        data should be summarised by the median.
  \item Standard deviation is variance in the original units, which is why it is
        the one usually reported.
  \item The $1.5\times$IQR rule identifies \emph{candidate} outliers. Whether to
        remove them is a domain decision, never a statistical one --- in security
        they are the target.
  \item Comparing mean with median is a free, instant diagnostic of skew.
  \item Normal, uniform, binomial, Poisson and long-tailed shapes each imply
        different valid summaries. Assuming normality by default is a common and
        costly error.
\end{tight}
\end{chaptersummary}

\begin{reviewq}
  \item \textbf{(MCQ)} Which measure of centre is unaffected by a single extreme
        value? (a)~Mean (b)~Median (c)~Variance (d)~Standard deviation.
  \item \textbf{(MCQ)} If mean $>$ median, the distribution is: (a)~left-skewed
        (b)~symmetric (c)~right-skewed (d)~uniform.
  \item \textbf{(Short)} State the $1.5\times$IQR rule and explain why its output
        should be called ``candidates'' rather than ``outliers''.
  \item \textbf{(Short)} Why is variance divided by $n-1$ rather than $n$ for a
        sample?
  \item \textbf{(Short)} Give one variable from IoT telemetry that is
        approximately normal and one that is long-tailed, and say how your
        reporting differs between them.
  \item \textbf{(Applied)} Session durations (minutes) for one user:
        $3, 4, 4, 5, 5, 6, 7, 9, 12, 140$. Compute Q1, Q3, IQR and the fences;
        identify flagged values; then argue both for keeping and for removing the
        140, naming the domain in which each argument wins.
  \item \textbf{(Applied)} A dashboard reports ``average time on task: 47
        minutes'' for an online course. Explain what could be wrong with this
        headline and propose two better numbers to display alongside it.
\end{reviewq}
